\documentclass[11pt,a4paper]{article}

\usepackage[utf8]{inputenc}

\usepackage[T1]{fontenc}

\usepackage[polish]{babel}

\usepackage{lmodern}

\usepackage{geometry}

\geometry{a4paper,margin=27mm}

\usepackage{microtype}

\usepackage[hidelinks]{hyperref}

% Portfolio: MSC-AR-P13; source: prospective ARC-XR research scope.

\title{Generatywna Neuroadaptive Art Therapy w środowisku XR}

\author{Mateusz Pomianek}

\date{Wersja 0.3.0, 21 listopada 2026}

\begin{document}

\maketitle

\begin{abstract}

Przedmiotem opracowania jest Generatywna Neuroadaptive Art Therapy w środowisku XR. Jak generatywne otoczenie wizualne reagujące na sygnały użytkownika wpływa na ocenę doświadczenia emocjonalnego? Porównany zostanie przebieg interakcji w generatywnym otoczeniu wizualnym dostosowującym obrazy do sygnałów użytkownika i w analogicznej scenie statycznej. Głównym punktem końcowym jest subiektywna ocena komfortu, zaangażowania i sprawczości podczas interakcji artystycznej. Opis ma charakter wstępny i nie przedstawia niezweryfikowanych wyników.

\end{abstract}

\section{Wprowadzenie}

Rozwój adaptacyjnych systemów XR stwarza możliwość badania interakcji człowieka z informacją przestrzenną, ale samo zastosowanie AR, VR lub modeli sztucznej inteligencji nie przesądza o przewadze poznawczej czy ergonomicznej. W rozpatrywanym zagadnieniu, Generatywna Neuroadaptive Art Therapy w środowisku XR, kluczowe jest rozdzielenie właściwości systemu od zmienności zachowania użytkownika. Porównanie należy zaprojektować tak, aby mierzyć jednoznaczne zjawisko, a nie tylko ogólną ocenę technologii.

\section{Cel pracy i pytanie badawcze}

Jak generatywne otoczenie wizualne reagujące na sygnały użytkownika wpływa na ocenę doświadczenia emocjonalnego? Projekt wymaga krytycznego przeglądu aktualnego stanu badań oraz zdefiniowania nowości w stosunku do innych tematów portfolio XR.

\section{Zakres i odrębny wkład naukowy}
Doświadczenie interaktywnej instalacji artystycznej reagującej na sygnały użytkownika. Termin terapia w nazwie koncepcji nie oznacza zatwierdzonej interwencji terapeutycznej. Wkład dotyczy poczucia wpływu w generatywnej sztuce XR, nie skuteczności psychoterapii.
\section{Hipotezy badawcze}
\paragraph{H1.} Kontyngentna wizualizacja generatywna zwiększa poczucie sprawczości w porównaniu ze sceną o identycznej estetyce bez reakcji.
\paragraph{H2.} Przewaga zależy od rzeczywistej kontyngencji sygnału i jest słabsza przy generowaniu niezależnym od zachowania.
\paragraph{H3.} Poprawa sprawczości nie musi powodować poprawy dobrostanu ani mieć efektu terapeutycznego.
\section{Projekt badania i plan eksperymentów}
\subsection{A. Operacjonalizacja i kwalifikacja danych}
Doświadczenie interaktywnej instalacji artystycznej reagującej na sygnały użytkownika. Termin terapia w nazwie koncepcji nie oznacza zatwierdzonej interwencji terapeutycznej. Jednostka wnioskowania: uczestnik × sesja estetyczno-interaktywna. W protokole zamrozić definicję i mianownik wyniku oraz źródło prawdy odniesienia, zgodność znaczników czasu, prawa do danych i zasady rekrutacji.
\subsection{B. Eksperyment porównawczy}
Stworzyć dwa warianty tego samego środowiska: reagujący na jawny sygnał i kontrolny odtwarzający te same obrazy niezależnie od użytkownika. Estymand pierwszorzędowy: różnica zdefiniowanej skali poczucia sprawczości między sceną generatywnie reagującą a statyczną. Przed testem zdefiniować randomizację, próbę, punkt końcowy i warunki wykluczenia bez selekcji wyniku post hoc.
\subsection{C. Kontrole i analiza odporności}
Nie wymagać EEG jeśli prosty sygnał sterowania wystarcza; oddzielić responsywność od atrakcyjności wizualnej. Zarejestrować preferencje artystyczne, dyskomfort, nudności i możliwość przerwania sesji. W rejestrze uwzględnić braki, wstrzymania, opóźnienia i niepełne próby zamiast usuwać je z analizy.
\subsection{D. Walidacja na niezależnym materiale}
Weryfikować na niezależnych sesjach lub osobach, bez dostępu modelu do danych końcowych. Model mieszany zwalidowanej skali sprawczości; zgłaszać negatywne odczucia i brak dowodów efektywności klinicznej. Dokumentować wyniki niesprzyjające hipotezie oraz rzeczywiste granice przeniesienia.
\section{Estymand, analiza statystyczna i decyzje}
różnica zdefiniowanej skali poczucia sprawczości między sceną generatywnie reagującą a statyczną. Jednostka: uczestnik × sesja estetyczno-interaktywna. Model mieszany zwalidowanej skali sprawczości; zgłaszać negatywne odczucia i brak dowodów efektywności klinicznej. H1 jest kandydatem hipotezy potwierdzającej, a H2/H3 pozostają eksploracyjne do formalnego planu wielokrotnych porównań. Liczebność uczestników, margines, próg i moc lub precyzja wymagają pilotażu oraz zamrożenia przed testem końcowym.
\section{Stan badań i granica nowości}
\begin{itemize}
\item XR personalization methods exist; a claim needs a defined multi-session or transfer contrast. \cite{xrprior1}.
\item AI adaptation in exposure contexts already exists; this plan is not a treatment efficacy study and requires dedicated clinical review. \cite{xrprior2}.
\end{itemize}
Wkład dotyczy poczucia wpływu w generatywnej sztuce XR, nie skuteczności psychoterapii. Odrębność od wcześniejszych publikacji i innych tematów programu należy sprawdzić przed rejestracją badania.
\section{Warunki bezpieczeństwa, sprzętu i zgodności}
Jest to plan prospektywny, nie dowód wykonania eksperymentu lub zatwierdzonej etyki. Udział osób, rejestrowanie biometrii i interwencje wymagają uprzedniej procedury etycznej i świadomej zgody. Nie wnioskować o emocjach ani zdrowiu z samych sygnałów. Standardowy Quest 3 nie ma okulometrii HMD, a Tobii Eye Tracker 5 służy do stanowiska ekranowego; każde połączenie zweryfikować przed protokołem.
\section{Materiał i metody}

Porównany zostanie przebieg interakcji w generatywnym otoczeniu wizualnym dostosowującym obrazy do sygnałów użytkownika i w analogicznej scenie statycznej. Wymagane jest zdefiniowanie jednostki analizy, zasad doboru przypadków, wersji scenariusza i źródeł pomiaru. Rozkład warunków oraz reguły wykluczenia prób należy określić przed pomiarem. W badaniach z udziałem ludzi niezbędna jest uprzednia właściwa ocena etyczna i świadoma zgoda. W testach wyłącznie symulacyjnych należy oddzielić prawdę symulacyjną od obserwacji empirycznej.

\subsection{Wynik główny i kryteria jakości}

Subiektywna ocena komfortu, zaangażowania i sprawczości podczas interakcji artystycznej. Pomiar powinien zachowywać jawne jednostki, znaczniki czasu i reguły oceny poprawności. Przy modelach predykcyjnych stosować rozłączne partycje oceny; przy eksperymentach interwencyjnych oddzielić skuteczność predykcji od przyczynowego wpływu interfejsu.

\section{Plan analizy}

Analiza kontrastów warunków i zmienności indywidualnej, w tym preferencji estetycznych. W wynikach należy raportować wielkość efektu lub błąd estymacji z przedziałem niepewności, liczbą obserwacji i odsetkiem odrzuconych prób. Analizy dodatkowe wymagają jasnego oznaczenia jako eksploracyjne. Nie wybierać progów decyzyjnych ani modelu na podstawie obserwacji ze zbioru testowego.

\section{Wyniki}

W obecnej wersji nie raportuje się wyników pomiarów, tabel statystycznych ani parametrów wyznaczonych z nieudokumentowanych danych. To kierunek badawczy, a nie zatwierdzony protokół badania ani zgłoszony artykuł.

\section{Dyskusja}

Interpretacja przyszłych wyników musi uwzględnić alternatywne mechanizmy i zależność wniosków od jakości danych, warunków zadania oraz ograniczeń sprzętu. Wartość potencjalnego rozwiązania zależy nie tylko od trafności modelu, lecz także od jego opóźnienia, stabilności i wpływu na zachowanie człowieka. Wyniki testów syntetycznych nie mogą być automatycznie uogólniane na populację uczestników.

\section{Ograniczenia badania}

Projekt nie stanowi walidacji psychoterapii ani leczenia; indywidualne preferencje silnie zakłócają wnioski. Granice nowości i trafności ekologicznej wymagają niezależnej weryfikacji. Żadna część szkicu nie dowodzi przeprowadzenia zbierania danych, uzyskania zgód etycznych ani gotowości manuskryptu do zgłoszenia.

\section{Wnioski}

Przedstawiono problem, kontrast i sposób jego planowanej oceny. Rozstrzygnięcie postawionego pytania wymaga przygotowania protokołu, zweryfikowanych źródeł oraz poprawnie udokumentowanych wyników. Bez tych materiałów nie można sformułować konkluzji empirycznych.

\begin{thebibliography}{9}
\bibitem{xrprior1} Adapting to the User: A Systematic Review of Personalized Interaction in VR (2025), arXiv:2510.13123. \url{https://arxiv.org/abs/2510.13123}.
\bibitem{xrprior2} Artificial intelligence (AI) for virtual reality exposure therapy (VRET): A systematic review (2026), Translational Psychiatry. \url{https://www.nature.com/articles/s41398-026-03936-4}.
\end{thebibliography}
\end{document}

